\documentclass[a4paper,12pt]{article}
\usepackage[T2A]{fontenc}
\usepackage[utf8]{inputenc}
\usepackage[russian]{babel}
\usepackage{amsmath,amssymb,mathtools}
\usepackage{geometry}
\geometry{left=22mm,right=20mm,top=22mm,bottom=24mm,headheight=26pt}
\usepackage{helvet}
\renewcommand{\familydefault}{\sfdefault}
\usepackage{microtype}
\usepackage{array,booktabs,tabularx}
\usepackage{enumitem}
\usepackage{setspace}
\usepackage{tikz}
\usepackage{xcolor}
\usepackage{hyperref}
\usepackage{fancyhdr}
\definecolor{accent}{HTML}{007F81}
\definecolor{deepaccent}{HTML}{005C60}
\definecolor{softgray}{HTML}{52585D}
\definecolor{errorred}{HTML}{A52D2D}
\hypersetup{colorlinks=true,linkcolor=deepaccent,pdfauthor={Лёвин Артём Александрович},pdftitle={Ревью домашней работы — 15 ноября 2026}}
\setlength{\parindent}{0pt}
\setlength{\parskip}{7pt plus 2pt minus 1pt}
\setstretch{1.11}
\setlist[enumerate]{leftmargin=*,itemsep=4pt,topsep=4pt}
\setlist[itemize]{leftmargin=*,itemsep=3pt,topsep=3pt}
\pagestyle{fancy}
\fancyhf{}
\fancyhead[L]{\small\textcolor{softgray}{Математика · профильный ЕГЭ}}
\fancyhead[R]{\small\textcolor{softgray}{15.11.2026}}
\fancyfoot[C]{\small\textcolor{softgray}{\thepage}}
\renewcommand{\headrulewidth}{0.3pt}
\renewcommand{\headrule}{\hbox to\headwidth{\color{accent}\leaders\hrule height \headrulewidth\hfill}}
\newcommand{\tasktitle}[2]{\phantomsection\label{task:#1}\section*{Задание №#1. #2}}
\newcommand{\err}{\textcolor{errorred}{\textbf{Ошибка:}}\ }
\newcommand{\idea}{\textcolor{deepaccent}{\textbf{Ключевая идея:}}\ }
\newcommand{\outcome}{\textcolor{deepaccent}{\textbf{Итог:}}\ }
\newcommand{\lab}[1]{\par\noindent\textbf{#1}\par\vspace{-4pt}}
\begin{document}
\thispagestyle{empty}
\vspace*{28mm}
\begin{center}
{\large\textcolor{deepaccent}{МАТЕМАТИКА}}\\[20mm]
{\LARGE Ревью домашней работы}\\[12mm]
{\Large\textcolor{accent}{Комплексное повторение}}\\[5mm]
{\normalsize Подготовка к ЕГЭ по профильной математике}
\vfill
{\normalsize Ксения Клыкова}\\[9mm]
{\normalsize 15 ноября 2026 г.}\\[19mm]
{\small Лёвин Артём Александрович}\\[1mm]
{\small эксклюзивно для Ксении Клыковой}\\[17mm]
\begin{tikzpicture}[x=1cm,y=1cm]
\draw[accent,line width=0.85pt] (0,0) .. controls (1.2,.5) and (2.2,-.3) .. (3.3,0)
  .. controls (4.4,.42) and (5.3,-.32) .. (6.4,0);
\end{tikzpicture}
\end{center}
\newpage

\section*{Сводная таблица}
В таблице приведены задания, для которых требуется исправление или завершение решения. Номера ведут к подробным разборам.
\par\vspace{4mm}
\noindent\begingroup
\small
\renewcommand{\arraystretch}{1.4}
\begin{tabularx}{\dimexpr\linewidth-3pt\relax}{@{}>{\centering\arraybackslash}p{13mm}>{\raggedright\arraybackslash}p{32mm}X>{\raggedright\arraybackslash}p{23mm}>{\raggedright\arraybackslash}p{27mm}@{}}
\toprule
№ & Тема & Суть замечания & Ваш ответ & Правильный ответ \\
\midrule
\hyperref[task:7]{7} & Вероятность & Подставлено значение $0{,}3$ вместо $0{,}21$. & $0{,}973$ & $0{,}990739$ \\
\hyperref[task:8]{8} & Случайное блуждание & Рассмотрен один шаг; вероятность достижения $-1$ за все последующие шаги не найдена. & Не указан & $0{,}3$ \\
\hyperref[task:10]{10} & Неравенство Чебышёва & Обозначен метод, но расчёт верхней оценки отсутствует. & Не указан & Не более $0{,}16$ \\
\bottomrule
\end{tabularx}
\endgroup\par
\vfill
\textcolor{softgray}{\small Для каждого задания приведены исправление, полное решение и ровно одно новое упражнение на закрепление.}
\newpage
\thispagestyle{plain}
\vspace*{0.32\textheight}
\begin{center}
{\LARGE\textcolor{deepaccent}{Разбор ошибок}}
\end{center}
\newpage

\tasktitle{7}{Вероятность хотя бы одного исправного элемента}
\lab{Текст задания}
Фонарь состоит из трёх ламп. Вероятность перегорания каждой лампы в течение года равна $0{,}21$. Найдите вероятность того, что хотя бы одна лампа за год не перегорит. В стандартной модели задачи события перегорания ламп считаются независимыми.

\lab{Ваше решение}
\[
P(\text{перегорят все})=0{,}3\cdot0{,}3\cdot0{,}3=0{,}027,
\]
\[
P(\text{хотя бы одна не перегорит})=1-0{,}027=0{,}973.
\]
\lab{Ваш ответ}
$0{,}973$.

\lab{Ошибка}
\err При вычислении вероятности перегорания всех ламп использовано $0{,}3$, хотя в условии дано $0{,}21$.

\lab{Где именно возникает ошибка}
\textbf{Последний правильный шаг:} выбор противоположного события. Чтобы найти вероятность того, что хотя бы одна лампа сохранится, удобно сначала рассчитать вероятность перегорания всех трёх.

\textbf{Первый неправильный шаг:}
\[
P(\text{перегорят все})=\underbrace{0{,}3\cdot0{,}3\cdot0{,}3}_{\text{неверное исходное число}}.
\]

\lab{Почему это неверно}
Число $0{,}21$ означает, что вероятность перегорания \emph{одной} лампы за год составляет $21\%$. Число $0{,}3$ соответствует уже $30\%$. Подмена исходных данных меняет итог, даже если все следующие действия выполняются правильно.

Правило умножения вероятностей независимых событий применено по смыслу верно. Нарушение произошло при переносе числа из условия.

\lab{Ключевая идея}
\idea Для события «хотя бы одна лампа не перегорит» противоположным будет событие «перегорят все три». Если $A$ --- искомое событие, то
\[
P(A)=1-P(\overline{A}).
\]

\lab{Исправление от последнего правильного шага}
Сохраняем выбранный способ, заменяя число на заданное:
\[
P(\overline A)=0{,}21\cdot0{,}21\cdot0{,}21.
\]

\lab{Полное правильное решение}
\begin{enumerate}
\item Вероятность перегорания одной лампы $q=0{,}21$.
\item По независимости вероятность перегорания всех трёх:
\[
q^3=(0{,}21)^3=0{,}009261.
\]
\item Вероятность противоположного события:
\[
P(A)=1-0{,}009261=\boxed{0{,}990739}.
\]
\end{enumerate}

\lab{Проверка}
Результат близок к единице: одновременное перегорание всех трёх ламп при вероятности $0{,}21$ для каждой --- сравнительно редкое событие. Проверка через дополнение: $0{,}990739+0{,}009261=1$.

\outcome Вероятность того, что хотя бы одна лампа не перегорит, равна $0{,}990739$.

\lab{Задача на закрепление}
Три независимые лампы работают в течение года. Вероятность перегорания каждой равна $0{,}2$. Найдите вероятность того, что за год хотя бы одна лампа не перегорит.
\newpage

\tasktitle{8}{Когда последовательность когда-нибудь достигнет заданного числа}
\lab{Текст задания}
Первый член случайной последовательности равен $0$. На каждом шаге значение увеличивается на единицу с вероятностью $p=\frac{10}{13}$ и уменьшается на единицу с вероятностью $1-p$. Найдите вероятность того, что \emph{какой-либо} член последовательности окажется равным $-1$.

\lab{Ваше решение}
Выписаны значения $0$, $1$, $-1$ и вероятности первого перехода: $\frac{10}{13}$ для шага вверх и $1-p$ для шага вниз. Итоговый ответ не получен.

\lab{Ваш ответ}
Не указан.

\lab{Ошибка}
\err Решение остановлено на вероятностях первого шага. В условии требуется вероятность достижения $-1$ \emph{когда-либо}, в том числе после нескольких шагов.

\lab{Где именно возникает ошибка}
\textbf{Последний правильный шаг:} указано, что после первого шага последовательность может попасть в $1$ с вероятностью $\frac{10}{13}$ или в $-1$ с вероятностью $1-p=\frac{3}{13}$.

\textbf{Первый отсутствующий шаг:} учёт путей, которые сначала уходят в положительную область, а затем всё-таки достигают $-1$. Достоверно ошибочного равенства в записи нет; решение осталось незавершённым.

\lab{Почему недостаточно одного шага}
Например, возможна последовательность перемещений
\[
0\longrightarrow1\longrightarrow0\longrightarrow-1.
\]
Она приводит к искомому значению за три шага. Её вероятность положительна:
\[
\frac{10}{13}\cdot\frac{3}{13}\cdot\frac{3}{13}>0.
\]
Поэтому вероятность встретить $-1$ когда-нибудь больше вероятности попасть туда сразу на первом шаге.

\lab{Ключевая идея}
\idea Нужно учитывать все возможные моменты первого достижения $-1$. Поскольку шаг вверх вероятнее шага вниз, вероятность такого достижения будет меньше единицы.

\lab{Исправление от последнего правильного шага}
Обозначим $q=1-p=\frac{3}{13}$ и $u$ --- вероятность когда-нибудь достичь $-1$, начиная с $0$. После первого шага:
\begin{itemize}
\item с вероятностью $q$ мы сразу попадаем в $-1$;
\item с вероятностью $p$ попадаем в $1$, после чего требуется сначала вернуться в $0$ и затем достичь $-1$.
\end{itemize}
Вероятность пройти два таких этапа равна $u^2$ (правила шагов не зависят от текущего уровня). Поэтому
\[
u=q+pu^2,
\qquad \frac{10}{13}u^2-u+\frac{3}{13}=0.
\]
После умножения на $13$:
\[
10u^2-13u+3=0,
\qquad (10u-3)(u-1)=0.
\]
Получились два формальных кандидата: $u=\frac{3}{10}$ и $u=1$. Чтобы выбрать верный, обоснуем вероятность достижения через конечную границу.

\lab{Полное правильное решение и выбор корня}
\begin{enumerate}
\item Рассмотрим вспомогательную игру: движение останавливается при достижении $-1$ или положительного уровня $N$. Пусть $u_k$ --- вероятность попасть в $-1$ раньше $N$, если начальное значение равно $k$.
\item Граничные условия: $u_{-1}=1$ и $u_N=0$. При $0\le k<N$:
\[
u_k=pu_{k+1}+qu_{k-1}.
\]
\item Переносим слагаемые и группируем разности:
\[
p(u_{k+1}-u_k)=q(u_k-u_{k-1}).
\]
Значит, соседние разности образуют геометрическую прогрессию со знаменателем
\[
r=\frac qp=\frac{3/13}{10/13}=\frac3{10}.
\]
\item Так как сумма разностей от $u_{-1}$ до $u_N$ равна $-1$, получаем
\[
u_N-u_{-1}=(u_0-u_{-1})(1+r+\dots+r^N)=-1.
\]
Отсюда с использованием суммы геометрической прогрессии
\[
u_0=1-\frac{1-r}{1-r^{N+1}}=\frac{r-r^{N+1}}{1-r^{N+1}}.
\]
\item Вспомогательную верхнюю границу $N$ можно выбирать всё дальше от нуля. Если путь когда-нибудь дойдёт до $-1$, то он сделает это за конечное число шагов, поэтому при достаточно большой границе $N$ такой путь будет учтён.

В формуле для $u_0$ содержится степень $r^{N+1}=(0{,}3)^{N+1}$. При увеличении $N$ эта степень становится сколь угодно малой: например, $0{,}3^3=0{,}027$, а $0{,}3^6=0{,}000729$. Поэтому вероятность достижения $-1$ когда-либо равна
\[
\boxed{u=r=\frac{3}{10}=0{,}3}.
\]
Здесь использованы только свойства степеней и геометрической прогрессии.
\end{enumerate}

\lab{Проверка}
Подстановка в найденное уравнение:
\[
\frac3{13}+\frac{10}{13}\left(\frac3{10}\right)^2
=\frac3{13}+\frac9{130}
=\frac{39}{130}=\frac3{10}.
\]
Также $0{,}3>\frac3{13}$, что согласуется с наличием путей длиной более одного шага. Корень $u=1$ для движения с более вероятным шагом вверх не подходит, что подтверждено расчётом через конечную границу.

\outcome Вероятность когда-нибудь получить член последовательности, равный $-1$, составляет $0{,}3$.

\lab{Задача на закрепление}
Последовательность начинается с $0$. На каждом шаге значение увеличивается на единицу с вероятностью $\frac34$ и уменьшается на единицу с вероятностью $\frac14$. Найдите вероятность того, что когда-нибудь последовательность достигнет $-1$.
\newpage

\section*{Теоретическая вставка}
{\large\textcolor{deepaccent}{Неравенство Чебышёва: от простого к формуле}}\par\vspace{2mm}

\lab{1. Зачем нужна такая оценка?}
Представьте, что фабрика выпускает упаковки с наполнителем. По данным фабрики,
\emph{средняя} масса равна $1500$~г. Это совсем не означает, что каждая
упаковка весит ровно $1500$~г: одна может оказаться легче, другая --- тяжелее.

Например, массы трёх упаковок равны $1480$, $1500$ и $1520$~г.
Их среднее значение:
\[
\frac{1480+1500+1520}{3}=1500\ \text{г}.
\]
Отклонения этих упаковок от среднего составляют $20$, $0$ и $20$~г
по величине. Нас интересует вопрос: \emph{как часто отклонение
оказывается большим, например больше $100$~г?}

Если известны только среднее значение и характеристика разброса,
точную вероятность обычно определить нельзя. Зато можно найти
\textbf{гарантированную верхнюю границу} этой вероятности.

\lab{2. Что такое среднее и стандартное отклонение?}
\textbf{Среднее значение} --- число, вокруг которого группируются
результаты. Его обозначают $\mu$ (читается «мю»).

\textbf{Стандартное отклонение} --- число, которое характеризует
разброс значений вокруг среднего. Его обозначают $\sigma$
(читается «сигма»). Чем меньше $\sigma$, тем в среднем плотнее
значения расположены около $\mu$.

Для сравнения: в наборах $98$, $100$, $102$ и $80$, $100$, $120$
среднее одинаковое --- $100$, но во втором наборе разброс заметно больше.

В задачах ЕГЭ стандартное отклонение часто \emph{уже дано}.
Вычислять его с нуля для применения неравенства Чебышёва не требуется.

\lab{3. Как понять слова «отличается больше чем на 100»?}
Если средняя масса $1500$~г, то отклонение больше $100$~г возникает,
когда масса \textbf{меньше $1400$~г или больше $1600$~г}:
\[
1500-100=1400,\qquad 1500+100=1600.
\]
Например, упаковка массой $1370$~г подходит под условие,
поскольку отличается от среднего на $130$~г. Упаковка массой
$1550$~г не подходит: отклонение всего $50$~г.

\begin{center}
\begin{tikzpicture}[x=1cm,y=1cm]
 \draw[->,softgray,line width=.7pt] (-6,0)--(6,0);
 \draw[line width=3pt,errorred!72] (-5.7,0)--(-2.6,0);
 \draw[line width=3pt,accent!72] (-2.6,0)--(2.6,0);
 \draw[line width=3pt,errorred!72] (2.6,0)--(5.7,0);
 \foreach \x in {-2.6,0,2.6} \draw[softgray] (\x,-.12)--(\x,.12);
 \node[below=5pt] at (-2.6,0) {$1400$};
 \node[below=5pt] at (0,0) {$1500$};
 \node[below=5pt] at (2.6,0) {$1600$};
 \node[above=7pt,text=errorred] at (-4.1,0) {далеко};
 \node[above=7pt,text=deepaccent] at (0,0) {не дальше 100 г};
 \node[above=7pt,text=errorred] at (4.1,0) {далеко};
\end{tikzpicture}
\end{center}

\newpage
\section*{Как работает неравенство Чебышёва}
\lab{4. Сначала разберём обозначения}
$X$ --- масса случайно выбранной упаковки; $P$ --- вероятность;
$\mu$ --- средняя масса; $\sigma$ --- стандартное отклонение;
$a>0$ --- выбранный порог отклонения (например, $100$~г).

Запись $|X-\mu|$ означает \textbf{расстояние от значения $X$ до среднего}.
Вертикальные черты --- это модуль: расстояние всегда неотрицательное.
Например,
\[
|1370-1500|=130,\qquad |1550-1500|=50.
\]
Таким образом, $P(|X-\mu|>a)$ --- вероятность того, что
случайное значение окажется дальше чем на $a$ от среднего.

\lab{5. Главное правило}
\begin{center}
{\large\textcolor{deepaccent}{\bfseries Неравенство Чебышёва}}
\end{center}
Для случайной величины с конечной дисперсией и любого $a>0$:
\[
\boxed{\displaystyle
P\bigl(|X-\mu|\ge a\bigr)\le \frac{\sigma^2}{a^2}.}
\]
Читаем словами: \emph{вероятность отклониться от среднего хотя бы
на $a$ не больше квадрата стандартного отклонения, делённого
на квадрат $a$.}

В условии задачи может быть сказано \emph{«более чем на $a$»}.
Тогда это правило тоже подходит:
\[
\boxed{\displaystyle
P\bigl(|X-\mu|>a\bigr)\le \frac{\sigma^2}{a^2}.}
\]
Почему? Если отклонение \emph{больше} $a$, оно обязательно и
\emph{не меньше} $a$. Второе событие включает первое.

\textbf{Очень важно:} знак $\le$ читается «меньше или равно».
Полученный результат --- \emph{максимально допустимая оценка сверху
по этому правилу}. Он не сообщает точную вероятность.
Например, результат «не более $0{,}16$» допускает как $0{,}04$,
так и $0{,}12$; оба числа меньше $0{,}16$.

\lab{6. Первый короткий пример}
Средняя масса товара $100$~г, стандартное отклонение $5$~г.
Какова верхняя оценка вероятности отклонения больше чем на $20$~г?

Выписываем $\mu=100$, $\sigma=5$, $a=20$ и применяем правило:
\[
P(|X-100|>20)\le\frac{5^2}{20^2}
=\frac{25}{400}=0{,}0625.
\]
То есть вероятность такого отклонения \textbf{не более $6{,}25\%$}.
Это верхняя оценка, даже если реальные массы товаров нам неизвестны.

\newpage
\section*{Практика применения: запомнить и проверить}
\lab{7. Второй пример: что означает ответ в процентах?}
Средний объём напитка равен $500$~мл, стандартное отклонение
$10$~мл. Оценим вероятность того, что объём отличается от
среднего больше чем на $25$~мл.
\[
P(|X-500|>25)\le\frac{10^2}{25^2}
=\frac{100}{625}=0{,}16=16\%.
\]
\textbf{Вывод:} вероятность слишком большого отклонения
\emph{не превосходит $16\%$}. Иными словами, вероятность
отклонения \emph{не больше $25$~мл} составляет \textbf{не менее $84\%$},
потому что $100\%-16\%=84\%$.

Здесь обе фразы дают полезную информацию:
\begin{itemize}
 \item «\emph{Не более $16\%$}» --- о слишком больших отклонениях;
 \item «\emph{Не менее $84\%$}» --- о значениях внутри выбранного диапазона.
\end{itemize}

\lab{8. Алгоритм для задачи ЕГЭ}
\begin{enumerate}[label=\arabic*.]
 \item Найдите в условии \textbf{среднее} $\mu$, \textbf{стандартное отклонение}
 $\sigma$ и порог отклонения $a$.
 \item Проверьте, что $\sigma$ и $a$ выражены в одинаковых единицах
 (например, оба числа --- в граммах).
 \item Возведите $\sigma$ и $a$ в квадрат.
 \item Разделите: $\dfrac{\sigma^2}{a^2}$.
 \item Запишите ответ со словами \textbf{«вероятность не более...»}.
 Если нужно, переведите десятичную дробь в проценты.
\end{enumerate}

\lab{9. Типичные ошибки}
\begin{itemize}
 \item \textbf{Забыли квадрат.} В формуле $\sigma^2/a^2$, а не $\sigma/a$.
 \item \textbf{Перепутали среднее и разброс.} Значение $\mu$ помогает
 понять, что значит отклонение; в дробь подставляются $\sigma$ и $a$.
 \item \textbf{Приравняли вероятность к оценке.} Правильно:
 $P\le0{,}16$; писать $P=0{,}16$ без дополнительных данных нельзя.
 \item \textbf{Получили число больше единицы.} Например, при $\sigma=20$
 и $a=10$ дробь равна $4$. Вероятность всегда не больше $1$,
 поэтому оценка $P\le4$ ничего нового не сообщает.
\end{itemize}

\lab{Перед разбором задания №10}
Здесь $\sigma=40$~г и $a=100$~г, поэтому верхняя оценка
равна $40^2/100^2=0{,}16$. Теперь разберём само решение по шагам.

\newpage

\tasktitle{10}{Верхняя оценка по неравенству Чебышёва}
\begingroup\small\setlength{\parskip}{4pt}
\lab{Текст задания}
Средняя масса наполнителя в упаковке равна $1500$~г, стандартное отклонение равно $40$~г. С помощью неравенства Чебышёва оцените сверху вероятность того, что масса случайно выбранной упаковки отличается от средней более чем на $100$~г.

\lab{Ваше решение}
Записано название метода: «Неравенство Чебышева?». Вычислений нет.

\lab{Ваш ответ}
Не указан.

\lab{Ошибка}
\err Требуемая оценка отсутствует: названа подходящая теорема, но не выписана формула и не подставлены данные.

\lab{Где именно возникает ошибка}
\textbf{Последний правильный шаг:} выбор неравенства Чебышёва как метода решения.

\textbf{Первый отсутствующий шаг:} переход от стандартного отклонения $\sigma=40$~г к дисперсии $\sigma^2=1600$~г$^2$ и применение формулы оценки. Неверного вычисления в записи нет.

\lab{Почему этот шаг необходим}
По условию известно стандартное отклонение. Неравенство Чебышёва использует \emph{дисперсию}, то есть квадрат стандартного отклонения. Кроме того, «верхняя оценка» обозначает гарантированную границу для вероятности; точное распределение масс неизвестно, поэтому точную вероятность по этим данным вычислить нельзя.

\lab{Ключевая идея}
\idea Если $M(X)=\mu$, стандартное отклонение равно $\sigma$, а $a>0$, то
\[
P\bigl(|X-\mu|\ge a\bigr)\le \frac{\sigma^2}{a^2}.
\]
Событие со строгим отклонением «более чем на $a$» содержится в событии «не менее чем на $a$», поэтому для него действует такая же верхняя оценка.

\lab{Исправление от последнего правильного шага}
Используем выбранное неравенство и данные задачи:
\[
\mu=1500\,\text{г},\quad \sigma=40\,\text{г},\quad a=100\,\text{г}.
\]

\lab{Полное правильное решение}
\begin{enumerate}
\item Обозначим массу случайной упаковки через $X$.
\item Дисперсия массы:
\[
D(X)=\sigma^2=40^2=1600\ \text{г}^2.
\]
\item Квадрат допустимого отклонения:
\[
a^2=100^2=10000\ \text{г}^2.
\]
\item Применим неравенство Чебышёва:
\[
\begin{aligned}
P\bigl(|X-1500|>100\bigr)
&\le P\bigl(|X-1500|\ge100\bigr)\\
&\le\frac{1600}{10000}=\boxed{0{,}16}.
\end{aligned}
\]
\end{enumerate}

\lab{Проверка}
Единицы измерения сокращаются: граммы в квадрате делятся на граммы в квадрате. Значение $0{,}16$ лежит между $0$ и $1$. Оценка означает «не более $16\%$»; сама вероятность может быть меньше.

\outcome Верхняя оценка вероятности: $0{,}16$.

\lab{Задача на закрепление}
Средняя масса упаковки равна $1800$~г, стандартное отклонение равно $25$~г. Используя неравенство Чебышёва, найдите верхнюю оценку вероятности отклонения массы от средней более чем на $100$~г.

\endgroup
\vfill
{\small\textcolor{softgray}{Лёвин Артём Александрович · эксклюзивно для Ксении Клыковой}}
\end{document}